% SHARD-ID: AQM-90-ALGEBRAIC-TORUS-QUANT2
% SHARD-TITLE: Algebraic Torus Quant2 Test Case
% SHARD-SUMMARY: Builds the split algebraic torus arithmetic-field Quant2 layer from closed finite-residue points.
% SHARD-SUMMARY: Computes Laurent differentials and records exactly how derivative predicates select Weyl-Heisenberg labels.
% SHARD-SUMMARY: Derives the toric-code-like Cayley-square edge complex from named torus units and expresses the result as Weyl constraints.
% SHARD-KEYWORDS: algebraic torus, Gm, Quant2, Weyl-Heisenberg, Laurent polynomials, dlog, toric code, Cayley complex

\section{Algebraic Torus Quant2 Test Case}
\label{sec:algebraic-torus-quant2}

\subsection{Status and Source Anchors}

This shard develops the first careful test case for the algebraic torus.
Stacks supplies spectra, \(D(f)\), localization, standard opens, and Kähler
differentials in \path{references/algebraic_geometry/stacks_project_algebra.tex},
lines 2972--3211, 33793--33872, and 34310--34340; affine morphisms and
anti-equivalence in \path{references/algebraic_geometry/stacks_project_schemes.tex},
lines 841--963; and finite closed-point residues by the AQM-28
Nullstellensatz locator.  The term \(\mathrm d\log\) is anchored by
\path{references/log_geometry/stacks_project_tag_0FMU.html}, lines 214--223.
The Weyl step is AQM-14, AQM-28, AQM-78, and convention \texttt{(ay)}.  CSS
constraints are AQM-08 and AQM-89.

\subsection{Lamport Dependency Skeleton}
\[
\begin{array}{rcl}
D1&:&k=\mathbb F_q,\quad A_T=k[x,x^{-1},y,y^{-1}],\quad T=\Spec A_T.\\
D2&:&T=D(xy)\subset\Spec k[x,y],\quad T(k)=(k^\times)^2.\\
D3&:&E_T(U)=\bigoplus_{z\in U\cap |T|_{\rm cl}}\kappa(z)^2,\quad
      \rho_U:\operatorname{WH}(E_T(U))\to PU(\mathcal C_U).\\
L1&:&\text{Nonzero Laurent monomials have infinite full closed support.}\\
D4&:&\mathrm d\log x=x^{-1}dx,\quad \mathrm d\log y=y^{-1}dy.\\
L2&:&\Omega^1_{A_T/k}=A_T\,\mathrm d\log x\oplus A_T\,\mathrm d\log y.\\
L3&:&d(x^ay^b)=x^ay^b(\bar a\,\mathrm d\log x+\bar b\,\mathrm d\log y).\\
T1&:&d\text{-flat monomial labels need fibre saturation on finite supports.}\\
D5&:&\alpha,\beta\in k^\times,\quad G=\langle\alpha\rangle\times\langle\beta\rangle.\\
T2&:&G\text{ gives a Cayley-square complex selecting edge Weyl labels.}
\end{array}
\]

\subsection{The Torus as a Scheme}

Let \(A_T=k[x,x^{-1},y,y^{-1}]\).  Equivalently, \(A_T=k[x,y]_{xy}\), so
\[
  T=\Spec A_T=D(xy)\subset\mathbb A^2_k.
\]
Thus a prime of \(T\) is a prime of \(k[x,y]\) not containing \(x\) or \(y\).
A \(k\)-rational point is a \(k\)-algebra map \(A_T\to k\), hence is determined
by nonzero images of \(x\) and \(y\).  Therefore
\[
  T(k)=(k^\times)^2.
\]
For \(a,b\in k^\times\), the corresponding maximal ideal is
\[
  \mathfrak m_{a,b}=(x-a,y-b)A_T,\qquad \kappa(\mathfrak m_{a,b})=k.
\]

The torus group law is not a slogan.  It is the affine morphism with
\[
  m^\#(x)=x\otimes x,\quad m^\#(y)=y\otimes y,\quad
  e^\#(x)=e^\#(y)=1,\quad i^\#(x)=x^{-1},\quad i^\#(y)=y^{-1}.
\]
The group axioms reduce, under affine anti-equivalence, to associativity,
unit, and inverse identities for multiplication of Laurent monomials.  This is
the precise sense in which \(T\) is an algebraic group object used here.

\subsection{Raw Closed-Point Quant2}

For an open \(U\subset T\), the raw arithmetic-field Quant2 label group is
\[
  E_T(U)=\bigoplus_{z\in U\cap |T|_{\rm cl}}\kappa(z)^2,
\]
with finite support, and the arithmetic quantum field is
\[
  \rho_U:\operatorname{WH}(E_T(U))\longrightarrow PU(\mathcal C_U).
\]
This is exactly the AQM-28 closed-point construction applied to \(A_T\).
The rational support \(S_{\rm rat}=T(k)=(k^\times)^2\) is a finite subtheory
with \((q-1)^2\) residue sites, each of residue field \(k\).  It is not the full torus theory, because the closed points of
\(T\) are infinite: for every monic irreducible \(\pi(x)\ne x\), the ideal
\[
  (\pi(x),y-1)A_T
\]
is a closed point with residue \(k[x]/(\pi)\), and AQM-29 records infinitely
many such \(\pi\).

\paragraph{Lemma \(L1\).}
Every Laurent monomial \(x^ay^b\), \(a,b\in\mathbb Z\), is nonzero at every
closed point of \(T\).  Hence its full closed-point profile has infinite
support.

\emph{Proof.}
The elements \(x\) and \(y\) are units in \(A_T\).  Their images in every
residue field \(\kappa(z)\) are therefore nonzero units, as are all Laurent
monomials \(x^ay^b\).  Since the closed-point set is infinite, the support is
infinite.

Thus Laurent monomials define formal full profiles.  They define actual Weyl
operators only after restriction to a finite closed support \(S\), for example
\(S_{\rm rat}\) or \(G\subset T(k)\).

\subsection{The Derivation on Laurent Monomials}

Let \(d:A_T\to\Omega^1_{A_T/k}\) be the universal \(k\)-derivation.  Define
\[
  \mathrm d\log x=x^{-1}dx,\qquad
  \mathrm d\log y=y^{-1}dy.
\]
For every \(A_T\)-module \(M\), a \(k\)-derivation \(D:A_T\to M\) is determined
freely by \(D(x)\) and \(D(y)\), because the inverse relation gives
\[
  D(x^{-1})=-x^{-2}D(x),\qquad D(y^{-1})=-y^{-2}D(y).
\]
Conversely, arbitrary \(u,v\in M\) define a derivation by the Laurent Leibniz
rule with \(D(x)=u\), \(D(y)=v\).  Hence
\[
  \Omega^1_{A_T/k}\cong A_T\,dx\oplus A_T\,dy
  =A_T\,\mathrm d\log x\oplus A_T\,\mathrm d\log y.
\]

\paragraph{Lemma \(L3\).}
For \(a,b\in\mathbb Z\),
\[
  d(x^ay^b)=x^ay^b(\bar a\,\mathrm d\log x+\bar b\,\mathrm d\log y),
\]
where \(\bar a,\bar b\in k\) are the images of the integers.

\emph{Proof.}
The formula follows from the Leibniz rule for positive exponents.  The inverse
relations above give the negative-exponent case.  Add the \(x\)- and
\(y\)-parts.

Since \(x^ay^b\) is a unit and \(\mathrm d\log x,\mathrm d\log y\) are a basis,
\[
  d(x^ay^b)=0\quad\Longleftrightarrow\quad
  a\equiv b\equiv0\pmod p.
\]
Thus the derivative is highly nontrivial on Laurent families: it kills
\(p\)-power monomials but not ordinary coordinate monomials.

\subsection{Derivative Selection Is Not Automatically Label-Intrinsic}

For a finite exponent set \(M\subset\mathbb Z^2\), set
\[
  U_M=\langle x^ay^b:(a,b)\in M\rangle_k\subset A_T.
\]
For a finite support \(S\subset |T|_{\rm cl}\), the residue label map is
\[
  \lambda_{S,M}:U_M^2\to E_T(S).
\]
The flat derivative predicate selects
\[
  C_{\rm flat}(S,M)=
  \lambda_{S,M}\{(q,p)\in U_M^2:dq=dp=0\}\subset E_T(S).
\]
The selected Weyl-Heisenberg elements are exactly
\[
  \operatorname{Sel}_{\rm flat}(S,M)=\pi^{-1}(C_{\rm flat}(S,M)),
\]
where \(\pi:\operatorname{WH}(E_T(S))\to E_T(S)\) is the label map.

The crucial warning is fibre saturation.  Over \(k=\mathbb F_5\), on the full
rational support \(S=T(k)\), the functions \(x\) and \(x^5\) have the same
residue profile because \(a^5=a\) for every \(a\in\mathbb F_5\).  But
\[
  d(x)=x\,\mathrm d\log x\ne0,\qquad
  d(x^5)=5x^5\,\mathrm d\log x=0.
\]
Therefore the same Weyl label can have a non-flat representative and a flat
representative if the family contains both \(x\) and \(x^5\).  By AQM-88, the
image definition still gives a selected subset, but it is not a
representative-free derivative predicate on labels unless the chosen family
and predicate are checked to be fibre-saturated.

\subsection{Cayley Squares from Torus Units}

Choose units \(\alpha,\beta\in k^\times\), with orders \(N_x,N_y\), and set
\[
  G=\langle\alpha\rangle\times\langle\beta\rangle\subset T(k).
\]
Write a point as \(g=(u,v)\).  Define
\[
  Xg=(\alpha u,v),\qquad Yg=(u,\beta v).
\]
The two translations commute because multiplication in \(k^\times\) commutes:
\[
  XYg=YXg=(\alpha u,\beta v).
\]
Now form the finite incidence sets
\[
  S_0=G,\qquad S_1=G\times\{x,y\},\qquad S_2=G.
\]
Over \(\mathbb F_2\), define
\[
  \partial_1(g,x)=g+Xg,\qquad
  \partial_1(g,y)=g+Yg,
\]
and
\[
  \partial_2(g)=(g,x)+(Xg,y)+(Yg,x)+(g,y).
\]
Then
\[
\begin{aligned}
\partial_1\partial_2(g)
&=g+Xg+Xg+XYg+Yg+YXg+g+Yg\\
&=XYg+YXg=0 .
\end{aligned}
\]
This is the torus group law showing up as a boundary identity.

Put qubit Weyl labels on the edge set \(S_1\):
\[
  E_\square(G)=\mathbb F_2^{S_1}{}_q\oplus\mathbb F_2^{S_1}{}_p.
\]
The CSS subgroup selected by the square complex is
\[
  L_\square(G)=
  (\operatorname{im}\partial_1^T\oplus0)+(0\oplus\operatorname{im}\partial_2)
  \subset E_\square(G).
\]
By AQM-08, \(\partial_1\partial_2=0\) is exactly the isotropy condition.  The
selected and rejected Weyl-Heisenberg elements are therefore
\[
  \operatorname{Sel}_\square(G)=\pi^{-1}(L_\square(G)),\qquad
  \operatorname{Rej}_\square(G)=
  \pi^{-1}(E_\square(G)\setminus L_\square(G)).
\]

\subsection{Worked Finite Examples}

For \(k=\mathbb F_3\), choose \(\alpha=\beta=2\).  Both have order \(2\).
Thus \(G=T(\mathbb F_3)\), with
\[
  |S_0|=4,\qquad |S_1|=8,\qquad |S_2|=4.
\]
This is the \(2\times2\) periodic square complex.  It is already a valid edge
Weyl-Heisenberg constraint system, though it is very small.

For \(k=\mathbb F_5\), choose \(\alpha=\beta=2\).  Since
\(2,2^2=4,2^3=3,2^4=1\), both units have order \(4\).  Hence
\[
  G=T(\mathbb F_5),\qquad |S_0|=16,\quad |S_1|=32,\quad |S_2|=16.
\]
The face at \(g=(1,1)\) is
\[
  ((1,1),x)+((2,1),y)+((1,2),x)+((1,1),y),
\]
whose boundary cancels because both paths end at \((2,2)\).  This has the
same size as the \(k=4\) toric instance validated in AQM-05 and AQM-06, now
with provenance from torus group law plus the named unit \(2\).

For \(k=\mathbb F_2\), \(k^\times=\{1\}\).  No nontrivial finite step is
available on \(T(k)\).  The rational-support Cayley-square route therefore
degenerates at \(\mathbb F_2\), even though qubit edge labels may still be
placed on a nontrivial torus coming from a larger field such as \(\mathbb F_5\).

\subsection{Conclusion}

The algebraic torus gives a sharper test case than the constant finite group
scheme.  Raw Quant2 on \(T\) is the closed-point residue Weyl field of
AQM-28.  The rational support \(T(k)\) is only a finite truncation.  The
derivation has canonical logarithmic coordinates and kills exactly the
\(p\)-power Laurent directions in the monomial basis, but derivative
constraints on finite rational Weyl labels are not automatically
representative-free.  Finally, after choosing two finite unit steps, the torus
group law produces a square chain complex whose output is precisely a selected
subgroup of the edge Weyl-Heisenberg group.  This is the rigorous toric-code
analogue to pursue next, not as a bare consequence of \(\Spec A_T\), but as
\[
  (T,\ x,y,\ \alpha,\beta,\ \mathbb F_2\text{-edge labels})
  \leadsto L_\square\subset E_\square
  \leadsto \pi^{-1}(L_\square)\subset\operatorname{WH}(E_\square).
\]
